% !TEX program = xelatex
% 金华十校2026年11月高三模拟考试 语文试题
\documentclass[11pt,a4paper]{ctexart}
\usepackage{ywexam}

\begin{document}

\examtitle{金华十校2026年11月高三模拟考试}{语文试题卷}
\examnote{注意：本卷共三大题，23小题，满分150分。考试时间150分钟。请按规定用笔将所有的答案写在答题纸上。}

\sechead{一、阅读（71分）}

\subsechead{（一）阅读Ⅰ（19分）}
阅读下面的文字，完成1～5题。

\material{材料一：}
\begin{materialtext}
“昆仑”在中华文化象征体系中，常被描绘为天地之轴、神人交界之所、不死药产地，具有强烈的神圣性与秩序感。它是多层次、多维度、高度凝练的精神符号，承载着中华民族的宇宙观、政治理念、伦理价值与生命理想，其流变传承考察需置于千百年来宏大的观念流动脉络中。

当前昆仑文化研究存在一个关键方法论困境。许多学者用后世成熟、抽象的哲学和天文学概念，反推上古先民原初、具象的思维，存在“时代错置”的风险。西汉扬雄的《太玄经》和宋代的《集韵》，分别是中华文明“数理推演”和“分类归纳”的理性化产物，用其注解“昆仑”的源头，就如同用高等数学去解释一位幼儿对“多”的最初感知，忽略了思维演进的漫长过程。我们应从更符合思维发展的角度，对“昆仑”观念的起源进行“由具象到抽象”的演进式解读。

“昆仑”并非一开始就指代“天”或“天道”的抽象概念，先民的思维最初是具象的。最能吸引他们注意力、与生存息息相关的具象事物是太阳。太阳作为万物生长的源泉，在远古先民心中的地位至高无上。最初的“昆仑”观念，很可能源于先民对太阳长期、持续的观察。太阳东升西落的自然现象，可能引发中华先民“死而复生”的神话想象，进而产生追求“永恒”“不死”的观念。

金文“昆”有“飞鸟负日”义，本义为“比日同行”。《说文》释“仑”为：“仑，思也。从人从册。”本义指条理、次序。“昆仑”连在一起，可能指太阳每天轮转运行的规律和秩序。远古先民对日月星辰、山脉河流有自然崇拜，于是认为有神鸟背负太阳而行，才导致东方日出和西方日落。这是一种朴素的宇宙观或天象观测。在先民的具象思维中，太阳每日东升西落，其运动方式与飞鸟最为相似，将太阳想象为神鸟（如“金鸟”），是极其自然的联想。

因此，“昆仑”的观念或许源于先民对太阳（飞鸟）运行轨迹的观察和神话想象，其核心是“太阳的循环”。在许多史前遗址中，也有相关线索，如河姆渡遗址出土的距今约7000年的“双鸟朝阳”象牙蝶形器，以同心圆太阳纹为核心，两侧双鸟昂首振翅、拱卫朝阳，堪称“日鸟合一”观念的早期定调之作。青海宗日遗址出土的距今约5000年的鸟（鹰）纹彩陶壶，壶口饰三角太阳纹，肩部环绕变形鸟纹，蕴含飞鸟绕日的意象。此外，陕西泉护村遗址出土的距今约5500年的仰韶文化庙底沟类型鸟纹彩陶盆“金乌负日图”，四川成都金沙遗址出土的距今约3000年的“太阳神鸟”金饰，以及汉画像石上的金乌负日纹，共同印证了“飞鸟负日”“日鸟合一”等信仰数千年的演进脉络。

随着先民活动范围扩大和地理认知深入，神话中的“太阳之站”需要现实地理载体。高山因其高耸入云、最接近太阳（飞鸟）的视觉特征，成为最佳选择。加之，在先民眼中，天需要柱子撑着，地上最高的山，就能沟通天地，成为天柱。于是，“昆仑”这个神话概念被“投射”或“锚定”到现实世界中的某座或某系列高山上。一些位置特殊、形体庞大的山川，开始被赋予人文含义，依不同标准被分类组合，构成象征意义。而日落西山、水天相接、长河落日等景象，或许正是古人猜想“河出昆仑”的最初源头。

“昆仑”是一个世代层累的文化概念：既是先民眼中那个每日东升西落的太阳，也是他们心中那个浑圆的天穹，更是远处那座巍峨神圣的高山。其内涵经历了从“太阳的循环规律”（具象观察）到“天的浑圆形态”（抽象模拟），再到“通天的神圣功能”（仪式物化）的完整演进过程。
\end{materialtext}
\begin{flushright}（摘编自鄂崇荣《昆仑文化的源流与传承》）\end{flushright}

\material{材料二：}
\begin{materialtext}
7月9日，“探秘山海经·何以昆仑”沉浸式艺术展在青海美术馆开幕。展览融合传统文物展陈与数字交互技术，深挖昆仑文化与《山海经》深厚底蕴，以创新形式推动中华优秀传统文化创造性转化、创新性发展。

展览锚定昆仑文化主线，融合静态展品与数字互动体验，并联动“西海2261·河湟文化大集”“昆仑文化论坛”等系列活动。展览兼顾科普性与趣味性，面向青少年普及昆仑文化，让厚重的本土文脉走近大众。

展览划分了多个功能区域，全方位展现山海昆仑文化魅力。历史科普区陈列昆仑神话图文、神兽造像，展出晋代郭璞注本《山海经》仿真件、尕日塘秦刻石相关资料，梳理古籍海外传播脉络。艺术展陈区展出西周遂公盨仿真拓片、古法竹编灯笼，展出基于明代朱耷笔意创作的穷奇、英招水墨作品与龙生九子造像，展示传统微缩景观制作技艺。数字互动区打造“天之四灵”光影空间、“字在山海”体感交互装置、九尾狐悬浮影像，用科技活化上古神话意象。
\end{materialtext}
\begin{flushright}（摘编自人民网《沉浸式艺术展邀你领略昆仑文化魅力》，2026年7月10日）\end{flushright}

\question{1}{下列对材料相关内容的理解和分析，不正确的一项是（3分）}
\begin{choices}
\item “昆仑”作为中华文化中的一个精神符号，承载着多重民族价值观念和民族心理。
\item 用西汉《太玄经》注解昆仑文化的起源，符合“由具象到抽象”的演进式解读原则。
\item 河姆渡遗址的“双鸟朝阳”象牙蝶形器，可被看作“日鸟合一”观念的早期佐证实物。
\item “探秘山海经·何以昆仑”展览通过数字交互技术，推动了昆仑文化的大众化传播。
\end{choices}

\question{2}{根据材料内容，下列说法不正确的一项是（3分）}
\begin{choices}
\item 先民最初将“昆仑”想象为太阳循环往复的“站点”，而非现实中某一座具体的高山。
\item 先民对自然崇拜的原始心理，引发了其对日升日落原因的联想，认为有神鸟负日而行。
\item “昆仑”这一神话概念被投射到高山，主要因为高山在先民心中能沟通天地，成为天柱。
\item 作为一个文化概念，“昆仑”内涵的演进过程，实际上也体现了先民认知观的发展变化。
\end{choices}

\question{3}{下列选项中，最适合作为论据来支撑材料一“昆仑观念源于先民对日鸟运行的观察”这一观点的一项是（3分）}
\begin{choices}
\item 山东大汶口文化遗址出土的陶尊上，刻有“日—火—山”组合的象形符号。
\item 《山海经》中记载“昆仑之丘，是实惟帝之下都”，将昆仑描述为天帝的居所。
\item 汉代《淮南子》记录了“昆仑之丘，或上倍之，是谓凉风之山”的层级地貌。
\item 清代金石学家收录的《昆仑石刻》残片，留存有秦代相关的刻字痕迹。
\end{choices}

\question{4}{简要分析材料一的论证思路。（4分）}

\par\vspace{2.2em}

\question{5}{青海美术馆拟在艺术展的数字互动区，新增一组面向青少年的“昆仑文化溯源”数字模拟活动。请结合两则材料相关内容，补写下面表格内容，为该活动建言献策。（6分）}

\begin{center}
\begin{tabular}{|c|p{2.6cm}|p{5.4cm}|p{4.8cm}|}
\hline
 & 活动名称 & 活动形式 & 活动理由 \\
\hline
透力 & “日鸟合一”数字模拟活动 & \xq{2} & 帮助青少年理解昆仑观念“日鸟合一”的完整脉络 \\
\hline
话型 & \xq{1} & 参观者操作数字模拟按钮，将神话中的太阳“驿站”逐步对应到现实中高原山脉 & \xq{3} \\
\hline
\end{tabular}
\end{center}

\subsechead{（二）阅读Ⅱ（16分）}
阅读下面的文字，完成6～9题。

\material{文本一}
\begin{center}{\bfseries 菊花（节选）}\par\vspace{0.15em}【美】约翰·斯坦贝克\end{center}

\begin{materialtext}
萨利纳斯山谷笼罩在深秋的薄光之中。高高的铁栅栏后面，伊丽莎·艾伦正在自家花园修剪菊花。她三十五岁，面容瘦削坚毅，眼神清亮。一身臃肿的园艺装束：男式黑帽、厚底靴子，印花裙外罩着带口袋的灯芯绒围裙，手上戴着结实的皮手套。\uline{\xq{1}她剪去菊花老茎，动作热切而用力。}

栅栏外，亨利·艾伦正站在拖拉机棚前和两个男人说话。亨利是伊丽莎的丈夫，他正在商谈卖掉那三十头三岁的小公牛。伊丽莎细心排查菊花苗周围的蚜虫、蜗牛，不让虫害靠近植株。亨利伏在保护花园的栅栏上，赞叹她的园艺本领：“你很有天赋，今年有些菊花直径足有二十五厘米。真希望你种苹果也能有这般本事。”

伊丽莎眼神锐利，颇为自得：“我或许也能种好苹果。\uline{\xq{2}我母亲就有一双庄稼人的手，什么栽下去都能活。}”

亨利告诉她，他把三十头三岁牛按预期价格卖掉了。“今天周六，下午我们进城下馆子，再看场电影庆祝。”

她在方形沙床里翻土，挖出平行小沟，修剪菊花秧苗，整齐码放。这时，一阵车轮与马蹄声传来。一辆破旧老式四轮马车沿着河岸小路驶来，帆布车顶像大篷车，由一匹老马和一头小灰驴拉动，一条瘦骨嶙峋的杂种狗跟在车下。帆布上歪歪扭扭写着：锅、平底锅、刀、剪刀、割草机。

马车拐进牧场小路，停在花园栅栏前。赶车人是个高大的中年男人，须发泛白，西装破旧沾满油污，手掌布满老茧裂纹。他向伊丽莎问路，又表明身份：自己走南闯北，沿途替人家修锅磨剪刀。

男人询问伊丽莎有没有器物需要修补，伊丽莎态度坚决，一口回绝。眼看没有生意，他的目光落到沙床的菊花上。“夫人，这是什么植物？”

伊丽莎脸上的怒气和抗拒一扫而光。“噢，这是菊花，个头很大，有白色的，有黄色的。我年年都种，这些菊花是这一带所有菊花里最大的。”

男人伏身在栅栏上。“听我说，我认识一位夫人，她就住在这条路上。她有世界上最漂亮的花园，花园里面有各式各样的奇花，就是没有菊花。我上次帮她修一个铜底洗衣盆时，难度很大，但我修得很好，她跟我说：‘万一你碰到品种好的菊花，拜托你给我带些种子。’她就是这么对我说的。”

听到这话，伊丽莎眼睛闪闪发亮。她摘下帽子，深色长发散落下来，热情地说：“直接移栽幼苗比播种更合适，我可以给她一些。”她跑到屋后取来红色花盆，徒手刨土，把菊花幼苗仔细栽进湿润沙土，认真叮嘱男人：幼苗一月生根，要分株移栽，七月离地二十厘米打顶，九月末就能绽放花朵。

跪在泥土里，伊丽莎神情炽热，向男人诉说“种植妙手”的奇妙感受：指尖和植物融为一体，万物仿佛听从指尖召唤，不会出错。她又向往起男人漂泊四方的旅途生活：“夜深的时候，星光落满周身，热烈璀璨。我从来没有经历过这样的日子。”

男人淡淡回应，旅途也会有挨饿受苦的时候。伊丽莎回过神，把花盆递给他，又拿来两个旧铝水壶，请他修补，付给他五十美分工钱。男人就地拿出工具敲打修复水壶。

“你就睡在马车里面吗？”伊丽莎问。

“风吹日晒都不怕，马车里安稳得很。”

伊丽莎满是羡慕：“一定棒极了。真希望女人也可以过这种生活。”

男人直言：“\uline{\xq{3}这种生活不适合女人。}”

她的上嘴唇微微向上抬起，露出牙齿。“你怎么知道？你怎么能确定？”她反问。

修好水壶，男人登上马车。伊丽莎反复叮嘱路上务必给花盆沙土浇水。马车缓缓掉头离开。伊丽莎扶着栅栏久久凝望远去的车队，嘴唇无声翕动，低声呢喃：“再见。\uline{\xq{4}远方满是光明，那是光来的方向。}”她的低语声把她自己吓了一跳，赶紧看看周围有没有人听到她说话的声音。只有狗听到了。伊丽莎转过身，匆忙跑进屋。

她跑到浴室里，脱下身上的脏衣服，接着，用浮石用力擦洗身体，皮肤被搓得泛红。而后她换上最好的内衣、袜子与漂亮衣裙，仔细梳好头发，描眉涂口红，认真装扮自己，内心满怀对远行、对挣脱牧场狭小生活的憧憬。

马蹄声传来，亨利回来了。伊丽莎端坐在门廊，望着河边雾霭里金黄的柳树枝条，久久出神。亨利看到盛装的她，由衷赞叹：“天哪，伊丽莎，你看起来美极了！”

“美？你觉得我很美？你说‘美’是什么意思？”

亨利继续迈步向前。“我不知道。我的意思是你看起来很不一样。很健壮，也很开心。”

“我很健壮？你说‘健壮’是什么意思？”

他无助地说，“你在玩游戏吧？拿我作消遣。你看起来健壮到可以把一只小牛犊放在膝盖上按死，开心得会跟吃西瓜似的再把小牛犊一口吃掉。”

有那么一刻，她不再那么端庄。“亨利！别说这种话。你根本不知道自己在说什么。”她又变回原来的端庄模样。“我很健壮，”她自夸道，“我以前从来不知道自己有多健壮。”

亨利说，“我去开车，你去把外套穿好。”小汽车在河边那条泥泞小路上颠簸，惊起鸟群，把兔子赶进灌木丛。两只鹤在柳树林上方用力拍打着翅翼，然后停歇在河床上。

伊丽莎远远地看见路的前方有一团小黑点。她知道那是什么。

他们经过时，她尽量不去看它，但眼睛不受她的控制。她悲戚地自言自语道：“他也许把它们都扔在了路上，这不费什么事，一点都不费事。但他会留下花盆。”她如此解释道：“他想留着花盆，所以不可能把它们丢在路边。”

汽车转了个弯，她看到前方的车队了。她转过身，面朝丈夫，这样汽车经过时她就不用看那辆小篷车和那一列怪模怪样的队伍。

一眨眼的工夫，一切都结束了，事情已成过去。她没有回头。

“亨利，”她问道，“我们吃晚饭时喝点酒，好吗？”

“当然可以啦。天哪！肯定会很开心。”

她无力地瘫在座位上，转过脸，不再看亨利。“能喝到酒我就已经很满足了。我们要喝很多酒。”她立起外套的领子，以免亨利看到她像个老妇人似的虚弱地哭泣。
\end{materialtext}
\begin{flushright}（有删改）\end{flushright}

\question{6}{下列对文本相关内容和艺术特色的分析鉴赏，不正确的一项是（3分）}
\begin{choices}
\item 小说开篇写伊丽莎的园艺装束，既表现她劳作时的干练，也暗示她身上带有男性化、力量化的气质。
\item 亨利赞叹伊丽莎的菊花种得好，又希望她能种好苹果，可见他虽认可妻子的园艺才能，但更看重实际收益。
\item 伊丽莎听了补锅匠的话后眼睛发亮，表明她已把补锅匠视为精神知己，并在他身上看到了丈夫所缺少的浪漫气质。
\item 小说结尾伊丽莎关于喝酒的谈论，表现她在幻想破灭后的虚弱、失落与痛苦，试图借酒掩饰自己的悲伤。
\end{choices}

\question{7}{关于文中画横线的语句，下列说法不正确的一项是（3分）}
\begin{choices}
\item 句\xq{1}，“热切而用力”描写修剪动作，既表现伊丽莎对园艺的投入，也暗示她内心被压抑的生命力在劳作中释放。
\item 句\xq{2}借母亲“庄稼人的手”自夸，表现伊丽莎对自己种植能力的自信，也暗含对亨利“种苹果”实用评价的不服气。
\item 句\xq{3}是补锅匠对伊丽莎向往流浪生活的直接否定，反映当时社会对女性角色的刻板定位，也刺痛了伊丽莎。
\item 句\xq{4}是伊丽莎对补锅匠远行生活的赞美，表明她已下定决心离开丈夫和牧场，去追寻“光来的方向”。
\end{choices}

\question{8}{这篇小说细腻展现伊丽莎从生出热烈期盼，逐步走向怅惘失落的心理转变。请结合文本，简析作者是如何呈现这一心理转变过程的。（4分）}

\par\vspace{2.2em}

\question{9}{某文学研读小组对这篇小说的主题有两种看法：有人认为小说赞扬女性的觉醒，也有人认为小说哀叹女性命运的不幸。你认同哪一种？结合文本分析。（6分）}

\par\vspace{2.2em}

\subsechead{（三）阅读Ⅲ（21分）}
阅读下面的文言文，完成10～14题。

\material{材料一：}
\begin{materialtext}
晋文公将与楚人战，召舅犯\xq{1}问之，曰：“吾将与楚人战，彼众我寡，为之奈何？”舅犯曰：“臣闻\zd{之}：繁礼君子，不厌忠信；战阵之间，不厌诈伪。君其诈之而已矣。”文公辞舅犯，因召雍季而问之，曰：“我将与楚人战，彼众我寡，为之奈何？”雍季对曰：“\uline{以诈遇民，偷取一时，后必无复。}”文公曰：“善。”辞雍季，以舅犯之谋与楚人战以败之。归而行爵，先雍季而\zd{后}舅犯。群臣曰：“城濮之事，舅犯谋也。夫用其言而后其身，可乎？”文公曰：“此非若所知也。夫舅犯言，一时之权也；雍季言，万世之利也。”仲尼闻之，曰：“文公之霸也宜哉！既知一时之权，又知万世之利。”

或曰：雍季之对，不\zd{当}文公之问。凡对问者，有因问小大缓急而对也。所问高大，而对以卑狭，则明主弗受也。今文公问“以少遇众”，而对曰“后必无复”，此非所以应也。且文公不知一时之权，又不知万世之利。战而胜，则国安而身定，兵强而威立，虽有后复，莫大于此，万世之利奚患不至？战而不胜，则国亡兵弱，身死名息，拔拂今日之死不及，安暇待万世之利？待万世之利，在今日之胜；今日之胜，在诈于敌；诈敌，万世之利而已。故曰：雍季之对，不当文公之问。且文公又不知舅犯之言。舅犯所谓“不厌诈伪”者，不谓诈其民，谓诈其敌也。\wave{敌者所伐之国也后虽无复何伤哉}？

舅犯曰“繁礼君子，不厌忠信”者，忠所以爱其下也，信所以不欺其民也。夫既以爱而不欺矣，言孰善于此？然必曰“出于诈伪”者，军旅之计也。舅犯前有善言，后有战胜，故舅犯有二功而后论，雍季无一焉而先赏。“文公之霸，不亦宜乎？”仲尼不知善赏也。
\end{materialtext}
\begin{flushright}（节选自《韩非子·难一》）\end{flushright}

\material{材料二：}
\begin{materialtext}
子玉释宋围，\zd{从}晋师。楚既陈，晋师退舍，军吏请曰：“以君避臣\xq{2}，辱也。且楚师老矣，必败。何故退？”子犯曰：“二三子忘在楚乎？偃也闻之：战斗，直为壮，曲为老。未报楚惠而抗宋\xq{3}，\uline{我曲楚直，其众莫不生气，不可谓老。}若我以君避臣，而不去，彼亦曲矣。”退三舍避楚。楚众欲止，子玉不肯，至于城濮，果战，楚众大败。君子曰：“善以德劝。”
\end{materialtext}
\begin{flushright}（节选自《国语·晋语四》）\end{flushright}

[注]\xq{1}舅犯，即狐偃。狐偃字子犯，因是晋文公重耳的舅舅，故史书常称其为“舅犯”。狐偃曾随重耳流亡各国，为重耳重要谋臣。\xq{2}以君避臣：指晋文公躲避楚臣子玉。\xq{3}未报楚惠而抗宋：重耳流亡时，曾受楚成王款待，承诺如遇战事，愿退避三舍以报楚恩。此时，晋为救宋而与楚国发生冲突。

\question{10}{材料一画波浪线的部分有三处需要断句，请用铅笔将答题卡上相应位置的答案标号涂黑，每涂对一处给1分，涂黑超过三处不给分。（3分）}

\duanju{敌者A所伐B之C国也D后E虽无F复G何H伤哉}

\question{11}{下列对材料中加点的词语及相关内容的解说，不正确的一项是（3分）}
\begin{choices}
\item 之，代词，代指下引句，与《孟子》中“人之有是四端也”的“之”用法不同。
\item 后，以……为后，与《岳阳楼记》中“后天下之乐而乐”的“后”字用法相同。
\item 当，切合，与《燕歌行（并序）》中“身当恩遇常轻敌”的“当”字意思不同。
\item 从，追随，与《鸿门宴》中“沛公旦日从百余骑来见项王”的“从”字意思相同。
\end{choices}

\question{12}{下列对材料有关内容的概述，不正确的一项是（3分）}
\begin{choices}
\item 材料一详写晋文公战前征询、赏后解释情形，刻画出晋文公谋定后动的特点。
\item 材料一中，孔子正向评价晋文公“善赏”，借此评价表达对舅犯用诈的否定。
\item 材料一通过辨析概念，阐释诈敌非诈民，借此批驳雍季“不应用诈”的言论。
\item 两则材料引入“或曰”“君子曰”的评论，均借评论表达对狐偃计谋的肯定。
\end{choices}

\question{13}{把材料中画横线的句子翻译成现代汉语。（8分）}

（1）以诈遇民，偷取一时，后必无复。

（2）我曲楚直，其众莫不生气，不可谓老。

\question{14}{两则材料认为城濮之战晋国取胜的原因分别是什么？两者归因不同，分别体现了法家、儒家的什么思想观点？（4分）}

\par\vspace{2.2em}

\subsechead{（四）阅读Ⅳ（9分）}
阅读下面这首唐诗，完成15～16题。

\begin{center}{\bfseries 喜外弟卢纶见宿}\par\vspace{0.15em}【唐】司空曙\end{center}

\begin{poem}
静夜四无邻，荒居旧业贫。雨中黄叶树，灯下白头人。\\
以我独沉久，愧君相见频。平生自有分，况是蔡家亲\textsuperscript{[注]}。
\end{poem}

[注]蔡家亲：表亲。

\question{15}{下列对这首诗的理解和赏析，不正确的一项是（3分）}
\begin{choices}
\item 诗题中的“喜”字为本诗诗眼，奠定了全诗“喜相逢”的感情基调。
\item 首联诗人自述生活艰困，为下文写顾惜、感念卢纶的深厚情谊张本。
\item 颔联采用意象对照的方式，营造特定氛围，对仗工整，画面感极强。
\item 末句“蔡家亲”呼应诗题“外弟”一词，表现诗人对亲戚之谊的珍视。
\end{choices}

\question{16}{尾联“分”字，有注解说“读fēn”，也有注解说“读fèn”，你认为应读何音？请简要阐释理由。（6分）}

\par\vspace{2.2em}

\subsechead{（五）名篇名句默写（6分）}

\question{17}{补写出下列句子中的空缺部分。（6分）}

（1）“树”和“木”在古典诗词中是两个不同的意象，“树”往往带有湿润、绵密之意，如本诗“雨中黄叶树”中的“树”，而“木”常常带有干燥、疏朗感觉，如黄庭坚《登快阁》“\blank[6.5em]，\blank[6.5em]”句中的“木”。

（2）“灯”在古代诗文中意涵丰富，可代指烟火温情，可象征落寞和孤独，也可借喻希望和指引，比如白居易《琵琶行》中的“\blank[6.5em]，\blank[6.5em]”。

（3）本诗颈联表达了诗人沉沦下僚的感慨，在古代诗文中，抒发类似感慨的名句层出不穷，比如：“\blank[6.5em]，\blank[6.5em]”。

\sechead{二、语言文字运用（19分）}
阅读下面的文字，完成18～22题。

近日，有文化名人发起了一场“抄书”活动，报名者之踊跃，出乎意料。有人埋头数十天，硬是把一整本书从头到尾誊写完毕。其实，抄书，抄的从来不只字面意思。

四千余字的《上林赋》，一笔一画写下来，要耗去五六个小时。有人反复描摩“荡荡乎八川分流”，满脑子都是某个人的影子，写完那一刻，心意也就到了，比发一长串微信郑重太多；有人失意后抄《庄子》，抄到“安时而处顺，哀乐不能入也”时，\blank[6.5em]，\blank[6.5em]，\blank[6.5em]，\blank[6.5em]；还有人在受挫后抄《岳阳楼记》，“不以物喜，不以己悲”，\uline{范仲淹的豁达像暖流渗进来，堵在心口的那团郁结松动了。}翻翻社交分享，书单千差万别，指向的却是同一种须要：爱慕时抄赋，迷茫时抄经，忧愁时抄诗词，想让自己硬气一点时，有人选择抄《孙子兵法》。抄什么，某种程度上是在回答“我现在缺什么？”

更多人把抄书当作一种独处仪式。现代人白天被消息轮番轰炸，晚上关掉手机，千愁万绪随着铺满的纸张一泻而尽。\wave{\xq{1}笔尖划过纸面的粗粝触感，\xq{2}每个字会慢下来，\xq{3}心也跟着定下来。\xq{4}没有消息弹窗，\xq{5}没有点赞提醒，\xq{6}总有笔尖与纸面的摩擦声，\xq{7}让人从碎片化的时间里抽离出来，\xq{8}获得一段完整的、属于自己的安静时光。}

\question{18}{文中第二段有两个含错别字的词语，请找出并改正。（2分）}

\par\vspace{2.2em}

\question{19}{依次填入文中横线空白处的语句，衔接最恰当的一组是（3分）}
\xq{1}撕心裂肺的痛

\xq{2}居然淡了几分

\xq{3}泪滴在纸面上

\xq{4}三遍之后

\xq{5}擦干了继续

\begin{choices}
\item \xq{1}\xq{3}\xq{5}\xq{4}\xq{2}
\item \xq{2}\xq{4}\xq{1}\xq{3}\xq{5}
\item \xq{3}\xq{5}\xq{4}\xq{1}\xq{2}
\item \xq{5}\xq{4}\xq{2}\xq{1}\xq{3}
\end{choices}

\question{20}{文中画横线的句子有“化虚为实”的表达效果，请结合文本分析该效果是怎样实现的。（4分）}

\par\vspace{2.2em}

\question{21}{文中画波浪线的部分有两处表述不当，请指出其序号并做修改，使语言准确流畅，逻辑严密，可少量增删词语，不改变原意。（4分）}

\par\vspace{2.2em}

\question{22}{汉语语义随时代发展不断更新，《辞海》中“抄”有“誊写，照录原文”的释义，未收录“抄书”义项。请结合文本内容与时代背景，为某词典新增一个“抄书”词条，拟写出两条释义，每条释义表述不超过15字，并分别说明该释义的确定依据。（6分）}

\par\vspace{2.2em}

\sechead{三、写作（60分）}

\question{23}{阅读下面的材料，根据要求写作（60分）}

激情是一种可贵的状态和品质，往往能最大限度地发挥创造潜能。在工作、学习或生活中，年青人总是激情满怀的，同时也容易激情过度。我们要利用好、维护好这种可贵的激情。

以上材料引发了你怎样的联想和思考？请写一篇文章。

要求：选准角度，确定立意，明确文体，自拟标题；不要套作，不得抄袭；不得泄露个人信息；不少于800字。

\end{document}
